\documentclass[UTF8,12pt]{ctexart}
\usepackage[a4paper,margin=24mm]{geometry}
\usepackage{amsmath,amssymb,bm,graphicx,adjustbox}
\newcommand{\fitmath}[1]{\begin{adjustbox}{max width=\linewidth}$\displaystyle #1$\end{adjustbox}}
\setlength{\parindent}{0pt}
\setlength{\parskip}{.7em}
\sloppy
\allowdisplaybreaks
\begin{document}
\section*{2001 年考研数学三 · 参考答案}
\subsection*{2001年真题参考答案}

\subsection*{一、填空题}

（1）\(\displaystyle -\dfrac{\displaystyle \alpha}{\displaystyle \beta}\)。（2）\(\displaystyle W_t=1.2W_{t-1}+2\)。（3）\(\displaystyle -3\)。（4）\(\displaystyle \dfrac{\displaystyle 1}{\displaystyle 12}\)。（5）\(\displaystyle F\)；\(\displaystyle (10,\allowbreak 5)\)。


\subsection*{二、选择题}

（1）B。（2）D。（3）C。（4）D。（5）A。


\subsection*{三、}

\(\displaystyle \dfrac{\displaystyle \mathrm du}{\displaystyle \mathrm dx}=\dfrac{\displaystyle \partial f}{\displaystyle \partial x}-\dfrac{\displaystyle y}{\displaystyle x}\dfrac{\displaystyle \partial f}{\displaystyle \partial y}+\left[1-\dfrac{\displaystyle e^x(x-z)}{\displaystyle \sin(x-z)}\right]\dfrac{\displaystyle \partial f}{\displaystyle \partial z}\)。


\subsection*{四、}

\(\displaystyle \dfrac{\displaystyle 1}{\displaystyle 2}\)。


\subsection*{五、}

\(\displaystyle -\dfrac{\displaystyle 2}{\displaystyle 3}\)。


\subsection*{六、}

（1）当 \(\displaystyle p=-\dfrac{\displaystyle 4}{\displaystyle 5}\)，\(\displaystyle q=3\) 时，\(\displaystyle S\) 最大。（2）最大值为 \(\displaystyle \dfrac{\displaystyle 225}{\displaystyle 32}\)。


\subsection*{七、}

证明略（可考虑函数 \(\displaystyle F(x)=xe^{1-x}f(x)\)，利用积分中值定理和罗尔定理）。


\subsection*{八、}

\(\displaystyle \sum_{n=1}^{\infty}f_n(x)=-e^x\ln(1-x)\)，\(\displaystyle -1\le x<1\)。


\subsection*{九、}

（1）\(\displaystyle a=-2\)。（2）\(\displaystyle Q=\begin{pmatrix}\dfrac{\displaystyle 1}{\displaystyle \sqrt2}&\dfrac{\displaystyle 1}{\displaystyle \sqrt6}&\dfrac{\displaystyle 1}{\displaystyle \sqrt3}\\0&-\dfrac{\displaystyle 2}{\displaystyle \sqrt6}&\dfrac{\displaystyle 1}{\displaystyle \sqrt3}\\-\dfrac{\displaystyle 1}{\displaystyle \sqrt2}&\dfrac{\displaystyle 1}{\displaystyle \sqrt6}&\dfrac{\displaystyle 1}{\displaystyle \sqrt3}\end{pmatrix}\)，\(\displaystyle Q^{\mathrm T}AQ=\begin{pmatrix}3&0&0\\0&-3&0\\0&0&0\end{pmatrix}\)。


\subsection*{十、}

（1）二次型 \(\displaystyle f(x_1,\allowbreak x_2,\allowbreak \cdots,\allowbreak x_n)\) 的矩阵形式为 \(\displaystyle f(X)=\dfrac{\displaystyle 1}{\displaystyle |A|}X^{\mathrm T}\operatorname{adj}(A)X\)。证明略。（2）因为 \(\displaystyle A\) 和 \(\displaystyle A^{-1}\) 合同，所以 \(\displaystyle g(X)\) 与 \(\displaystyle f(X)\) 有相同的规范形。


\subsection*{十一、}

最多可装98箱。


\subsection*{十二、}

\(\displaystyle p(u)=\begin{cases}\dfrac{\displaystyle 1}{\displaystyle 2}(2-u),\allowbreak &0<u<2,\allowbreak \\0,\allowbreak &\text{其他}.\end{cases}\)

\end{document}
